\documentclass[border=9pt]{standalone}
\input{figure_preamble.tex}

\begin{document}
\begin{tikzpicture}

\def\ssize{1}
\def\boxsize{1.3}
\def\smallsize{0.13}

\tikzstyle{snode} = [fig/module,minimum width=0.28cm,minimum height=0.28cm,inner sep=0pt]

% =========================
% (a) ODE
% =========================
\begin{scope}
\draw [-{Latex[scale=1.2]},thick] (0,0) -- (13.5,0);
\draw [-{Latex[scale=1.2]},thick] (0,-1.5) -- (0,2.5);

\node [anchor=north east,align=right] at (-0.1,2.5) {\footnotesize{$\mathbf{x}(t)$}};
\node [anchor=north,align=left] at (13.3,-0.1) {\footnotesize{时间 $t$}};

% ODE 轨迹
\draw [Round Cap-Round Cap,blue,very thick,domain=0:12.7,samples=120]
plot (\x,{0.3-0.001*(((\x-1)/1.15)^4) * (\x-1) + 0.029375*(((\x-1)/1.15)^4) + 1/(1.15)^4*(- (323/1200)*((\x-1)^3) + 0.75*((\x-1)^2))});

% 图例
\node [anchor=west] (legend1) at (1.5,2.4) {\footnotesize{确定性轨迹}};
\draw [-,blue,very thick] ([xshift=-0.2cm]legend1.west) -- ([xshift=-0.8cm]legend1.west);

% 标题
\node [fig/caption,anchor=north] at (6.5,-2.0) {\small{(a) ODE（确定性轨迹）}};
\end{scope}

% =========================
% (b) SDE
% =========================
\begin{scope}[yshift=-6.5cm]
\draw [-{Latex[scale=1.2]},thick] (0,0) -- (13.5,0);
\draw [-{Latex[scale=1.2]},thick] (0,-1.5) -- (0,2.5);

\node [anchor=north east,align=right] at (-0.1,2.5) {\footnotesize{$\mathbf{x}(t)$}};
\node [anchor=north,align=left] at (13.3,-0.1) {\footnotesize{时间 $t$}};

% 定义中间轨迹函数
\def\traj#1{(0.3-0.001*(((#1-1)/1.15)^4)*(#1-1) + 0.029375*(((#1-1)/1.15)^4) + 1/(1.15)^4*(-(323/1200)*((#1-1)^3) + 0.75*((#1-1)^2))) * ((-#1/12.7)+1)}

% 定义噪声幅度函数（随时间逐渐增大）
\def\noise#1{1*(#1/12.7)}

% 先画噪声阴影带
\fill[orange!20,opacity=0.8]
  plot[domain=0:12.7,samples=120] (\x,{\traj{\x} + \noise{\x}})
  --
  plot[domain=12.7:0,samples=120] (\x,{\traj{\x} - \noise{\x}})
  -- cycle;

% 画上下边界虚线
\draw [dotted,gray!70,domain=0:12.7,samples=120]
  plot (\x,{\traj{\x} + \noise{\x}});
\draw [dotted,gray!70,domain=0:12.7,samples=120]
  plot (\x,{\traj{\x} - \noise{\x}});

% 画确定性轨迹
\draw [Round Cap-Round Cap,blue,very thick,domain=0:12.7,samples=120]
  plot (\x,{\traj{\x}});

% 噪声幅度标注箭头
\foreach \x in {10}{
    \pgfmathsetmacro{\p}{\traj{\x} + \noise{\x}}
    \pgfmathsetmacro{\q}{\traj{\x} - \noise{\x}}
    \node [circle,minimum size=2pt,inner sep=0,fill=blue] (pointp) at (\x,\p) {};
    \node [circle,minimum size=2pt,inner sep=0,fill=blue] (pointq) at (\x,\q) {};
    \draw [<->,blue] ([yshift=-1pt]pointp.south) -- ([yshift=1pt]pointq.north);
}

% 图例1：确定性轨迹
\node [anchor=west] (legend1) at (1.5,2.4) {\footnotesize{确定性轨迹}};
\draw [-,blue,very thick] ([xshift=-0.2cm]legend1.west) -- ([xshift=-0.8cm]legend1.west);

% 图例2：随机噪声
\node [anchor=west] (legend2) at (1.5,1.9) {\footnotesize{随机噪声}};
\node [anchor=east,inner sep=0,draw=gray!70,fill=orange!20,dotted,
       minimum width=0.6cm,minimum height=0.2cm]
       at ([xshift=-0.2cm]legend2.west) {};

% 噪声幅度文字
\node [anchor=south east] at (10.7,0.6) {\footnotesize{噪声幅度}};

% 标题
\node [fig/caption,anchor=north] at (6.5,-2.0) {\small{(b) SDE（含随机噪声）}};
\end{scope}

\end{tikzpicture}
\end{document}



% \documentclass[border=9pt]{standalone}
% \input{figure_preamble.tex}
% \begin{document}
% \begin{tikzpicture}

% \def\ssize{1}
% \def\boxsize{1.3}
% \def\smallsize{0.13}

% \tikzstyle{snode} = [fig/module,minimum width=0.28cm,minimum height=0.28cm,inner sep=0pt]

% \begin{scope}
% \draw [-{Latex[scale=1.2]},thick] (0,0) -- (13.5,0);
% \draw [-{Latex[scale=1.2]},thick] (0,-1.5) -- (0,2.5);
% \node [anchor=north east,align=right] (ylabel) at (0-0.1,2.5) {\footnotesize{$\mathbf{x}(t)$}};
% \node [anchor=north,align=left] (xlabel) at (13.3,-0.1) {\footnotesize{时间 $t$}};

% \draw [Round Cap-Round Cap,blue,very thick,domain=0:12.7,samples=100] plot (\x,{0.3-0.001*(((\x-1)/1.15)^4) * (\x-1) + 0.029375*(((\x-1)/1.15)^4) + 1/(1.15)^4*(- (323/1200)*((\x-1)^3) + 0.75*((\x-1)^2))});

% \node [anchor=west] (legend1) at (1.5,2.4) {\footnotesize{确定性轨迹}};
% \draw [-,blue,very thick] ([xshift=-0.2cm]legend1.west) -- ([xshift=-0.8cm]legend1.west);

% \node [fig/caption,anchor=north] (caption) at (6.5,-2.0) {\small{(a) ODE（确定性轨迹）}};

% \end{scope}

% \begin{scope}[yshift=-6.5cm]
% \draw [-{Latex[scale=1.2]},thick] (0,0) -- (13.5,0);
% \draw [-{Latex[scale=1.2]},thick] (0,-1.5) -- (0,2.5);
% \node [anchor=north east,align=right] (ylabel) at (0-0.1,2.5) {\footnotesize{$\mathbf{x}(t)$}};
% \node [anchor=north,align=left] (xlabel) at (13.3,-0.1) {\footnotesize{时间 $t$}};

% \draw [Round Cap-Round Cap,blue,very thick,domain=0:12.7,samples=100] plot (\x,{(0.3-0.001*(((\x-1)/1.15)^4) * (\x-1) + 0.029375*(((\x-1)/1.15)^4) + 1/(1.15)^4*(- (323/1200)*((\x-1)^3) + 0.75*((\x-1)^2))) * ((-\x/12.7) + 1)});

% \begin{pgfonlayer}{background}
% \draw [Round Cap-Round Cap,name path=linetop,dotted,domain=0:12.7,samples=100] plot (\x,{(0.3-0.001*(((\x-1)/1.15)^4) * (\x-1) + 0.029375*(((\x-1)/1.15)^4) + 1/(1.15)^4*(- (323/1200)*((\x-1)^3) + 0.75*((\x-1)^2))) * ((-\x/12.7) + 1) + 1 * (\x/12.7) });
% \draw [Round Cap-Round Cap,name path=linebottom,dotted,domain=0:12.7,samples=100] plot (\x,{(0.3-0.001*(((\x-1)/1.15)^4) * (\x-1) + 0.029375*(((\x-1)/1.15)^4) + 1/(1.15)^4*(- (323/1200)*((\x-1)^3) + 0.75*((\x-1)^2))) * ((-\x/12.7) + 1) - 1 * (\x/12.7) });
%     \tikzfillbetween[of=linetop and linebottom,on layer=background]{orange!20}
% \end{pgfonlayer}

% \foreach \x in {10}{
% \pgfmathsetmacro{\p}{(0.3-0.001*(((\x-1)/1.15)^4) * (\x-1) + 0.029375*(((\x-1)/1.15)^4) + 1/(1.15)^4*(- (323/1200)*((\x-1)^3) + 0.75*((\x-1)^2))) * ((-\x/12.7) + 1) + 1 * (\x/12.7) }
% \pgfmathsetmacro{\q}{(0.3-0.001*(((\x-1)/1.15)^4) * (\x-1) + 0.029375*(((\x-1)/1.15)^4) + 1/(1.15)^4*(- (323/1200)*((\x-1)^3) + 0.75*((\x-1)^2))) * ((-\x/12.7) + 1) - 1 * (\x/12.7) }
% \node [circle,minimum size=2pt,inner sep=0,fill=blue] (pointp) at (\x,\p) {};
% \node [circle,minimum size=2pt,inner sep=0,fill=blue] (pointq) at (\x,\q) {};
% \draw [<->,blue] ([yshift=-1pt]pointp.south) -- ([yshift=1pt]pointq.north);
% }

% \node [anchor=west] (legend1) at (1.5,2.4) {\footnotesize{确定性轨迹}};
% \draw [-,blue,very thick] ([xshift=-0.2cm]legend1.west) -- ([xshift=-0.8cm]legend1.west);

% \node [anchor=west] (legend2) at (1.5,1.9) {\footnotesize{随机噪声}};
% \node [anchor=east,inner sep=0,draw,fill=orange!20,thin,dotted,minimum width=0.6cm,minimum height=0.2cm] (legend2label) at ([xshift=-0.2cm]legend2.west) {};

% \node [anchor=south east] (noiselabel) at (10.7,0.6) {\footnotesize{噪声幅度}};

% \node [fig/caption,anchor=north] (caption) at (6.5,-2.0) {\small{(b) SDE（含随机噪声）}};

% \end{scope}

% \end{tikzpicture}
% \end{document}
% figure/figure-forward-sde-example.tex

% 前向 SDE 将双峰非高斯分布逐渐变换为标准高斯分布。
% 彩色曲线表示随机轨迹，灰色曲线表示确定性的概率流 ODE 轨迹。

% figure/figure-forward-sde-example.tex

% 说明：
% 1. 左侧为双峰非高斯分布 p_0(x)
% 2. 右侧为单峰高斯先验 p_T(x)
% 3. 灰色曲线表示光滑的 probability flow ODE
% 4. 彩色曲线表示更粗糙、更曲折的 SDE 样本轨迹

% 使用方式：
% \begin{figure}[!t]
%   \centering
%   \AMLFigure{figure-forward-sde-example}
%   \caption{...}
%   \label{fig:forward-sde-example}
% \end{figure}

